\section{Lecture 10 (February 5, 2025): Reconstruction of an affine group}\label{lec10}
\noindent\textit{Scribe: Jackson Morris}
\label{lec?(2-5)}

We prove the reconstruction statement announced in Lecture~\ref{lec9}.
\begin{theorem}\label{thm:group-reconstruction}
For an affine group scheme $G$ over a field $k$, the natural morphism
\[
 G\longrightarrow\underline{\operatorname{Aut}}^{\otimes}(\omega),
 \qquad \omega:\operatorname{Rep}_k(G)\to\mathrm{Vec}_k,
\]
is an isomorphism of group-valued functors. In particular, the tensor automorphism functor is representable.
\end{theorem}

\subsection*{Chevalley's theorem}
\begin{theorem}[Chevalley]
Every closed subgroup scheme $H$ of an affine algebraic group $K$ is the
stabilizer of a one-dimensional subspace in a finite-dimensional
representation of $K$.
\end{theorem}
See \cite[Chapter VIII, Theorem 13.1]{milneAGS}.

\begin{proof}[Proof of Theorem~\ref{thm:group-reconstruction}]
Fix $X\in\operatorname{Rep}_k(G)$ and put $H=G_X$ and
$K=\mathrm{GL}(X)$. As in Lecture~\ref{lec9}, the tensor subcategory
$\mathcal C_X$ generated by $X$ identifies with $\operatorname{Rep}(H)$,
and for every commutative $k$-algebra $R$ we have inclusions
\[
 H(R)\subseteq
 \underline{\operatorname{Aut}}^{\otimes}(\omega|_{\mathcal C_X})(R)
 \subseteq K(R).
\]
Let $\lambda$ be a tensor automorphism in the middle group and write
$g=\lambda_X\in K(R)$. By Chevalley's theorem, choose a
$K$-representation $W$ and a line $L\subseteq W$ with
$H=\operatorname{Stab}_K(L)$.

Every finite-dimensional $K$-representation is a subquotient of a finite
direct sum of tensor products of $X$ and $X^*$
\cite[Chapter VIII, Theorem 11.7]{milneAGS}. Thus $W|_H$ belongs to
$\mathcal C_X$, and the action of $g$ on $W\otimes R$ agrees with
$\lambda_W$: both are obtained from $g$ by duals, tensor products,
direct sums and passage to subquotients.

The line $L$ is an $H$-subrepresentation of $W$. Naturality for its
inclusion $i:L\hookrightarrow W$ gives
\[
 \begin{tikzcd}
 L\otimes R \arrow[r,"\lambda_L"] \arrow[d,"i\otimes1"'] &
 L\otimes R \arrow[d,"i\otimes1"]\\
 W\otimes R \arrow[r,"\lambda_W"'] & W\otimes R.
 \end{tikzcd}
\]
Since $\lambda_L$ is invertible, $g(L\otimes R)=L\otimes R$.
Hence $g\in H(R)$. Its action on $X$ is $\lambda_X$, which determines
all components on $\mathcal C_X$, so its action is exactly $\lambda$.
This proves
\[
 G_X\simeq\underline{\operatorname{Aut}}^{\otimes}
 (\omega|_{\mathcal C_X}).
\]
These isomorphisms commute with restriction when
$\mathcal C_X\subseteq\mathcal C_{X\oplus Y}$. The Hopf algebras
$\mathcal O(G_X)$ are finitely generated as $k$-algebras and exhaust
$\mathcal O(G)$, as proved in Lecture~\ref{lec9}. Passing to the
compatible inverse limits gives
\[
 G\simeq\varprojlim_X G_X
 \simeq\varprojlim_X
 \underline{\operatorname{Aut}}^{\otimes}(\omega|_{\mathcal C_X})
 \simeq\underline{\operatorname{Aut}}^{\otimes}(\omega).
\]
All constructions are natural in $R$ and respect the group operations.
\end{proof}

\subsection*{Functoriality}
A group homomorphism $f:G\to G'$ gives restriction of representations
\[
 f^*:\operatorname{Rep}(G')\longrightarrow\operatorname{Rep}(G),
\]
with the canonical identification $\omega_Gf^*=\omega_{G'}$.
\begin{corollary}
Group homomorphisms $G\to G'$ correspond to symmetric tensor functors $F:\operatorname{Rep}(G')\to\operatorname{Rep}(G)$ equipped with a specified tensor isomorphism $\alpha:\omega_GF\simeq\omega_{G'}$, up to isomorphisms respecting $\alpha$.
\end{corollary}
\begin{proof}
For each $R$, an element $g\in G(R)$ acts on every $F(V)_R$. Transport this action through $\alpha$ to obtain a tensor automorphism of $\omega_{G',R}$. Theorem~\ref{thm:group-reconstruction} identifies it with a unique element of $G'(R)$. Naturality and multiplication give a group-scheme map $f:G\to G'$, and the same construction shows that $\alpha$ intertwines $F$ with $f^*$.
\end{proof}

When $k$ is algebraically closed, every symmetric tensor functor
$F:\operatorname{Rep}(G')\to\operatorname{Rep}(G)$ is isomorphic to
restriction along a group homomorphism $G\to G'$. Indeed, the two fiber
functors $\omega_GF$ and $\omega_{G'}$ are tensor-isomorphic
\cite[Corollary 1.2]{WibmerTorsors}; choosing such an isomorphism and applying
the preceding corollary gives the homomorphism.

\subsection*{The general reconstruction problem}
We next prove the following extension, in Lectures~\ref{lec11}--\ref{lec13}.
\begin{theorem}[Neutral Tannaka reconstruction]\label{thm:neutral-tannaka}
Let $\mathcal C$ be an essentially small symmetric tensor category over $k$, and let $\omega:\mathcal C\to\mathrm{Vec}_k$ be an exact faithful $k$-linear strong symmetric monoidal functor to finite-dimensional vector spaces. Then
\[
 G=\underline{\operatorname{Aut}}^{\otimes}(\omega)
\]
is an affine group scheme, and $\omega$ lifts to a symmetric tensor equivalence $\widetilde\omega:\mathcal C\simeq\operatorname{Rep}_k(G)$ whose composition with the forgetful functor is $\omega$.
\end{theorem}
We first reconstruct a coalgebra from the abelian category and $\omega$, then use the tensor product, symmetry and duals to recover multiplication, commutativity and the antipode.
